\documentclass[11pt,a4paper]{coverletter}

\senderblock
  {Radheem Bin Razi}
  {Distributed Systems \& Cloud-Native Engineer}
  {Ilmenau, Germany \textbar\ \href{mailto:sheikh.radheem@gmail.com}{sheikh.radheem@gmail.com} \textbar\ \href{https://radheem.github.io/my-cv}{portfolio}}

\begin{document}

%% ===== ENGLISH =====
\recipient{ilert GmbH}{Hiring Team}{}

\opening{Dear Hiring Team,}

\vspace{8pt}\noindent\textbf{1. Why ilert GmbH?}\par\vspace{3pt}

ilert's shift toward an AI-first autonomous layer for incident response presents a compelling opportunity to apply my distributed systems background. I am eager to help your team scale a platform where AI actively investigates issues, reducing manual toil for DevOps engineers. I bring React, TypeScript, and Node.js experience to bridge frontend and backend concerns, while my work with Kafka and distributed pipelines aligns with your data architecture. My builder mentality is demonstrated by projects like IRS, where I automated workflows using FastMCP and build  information retrival system with data integrity, which required rigorous system design.

\vspace{8pt}\noindent\textbf{2. Why me?}\par\vspace{3pt}

\bullets

  \item Owned the IRS platform end-to-end: built Go microservices with gRPC, built an event driven system, and deployed on Kubernetes with Cilium.

  \item Integrated multiple third-party vendors via an adaptor pattern, ensuring reliable data flow between external vendors and the core information engine.

  \item Used Hatchet for managing workflows and webhooks for automation,

\bulletsend

I am currently studying at TU Ilmenau, my coursework is complete and I am starting my thesis allowing me full availability and immediate relocation to Cologne for the hybrid work.

\closing{Sincerely,}{Radheem Bin Razi}

\clearpage
\selectlanguage{ngerman}

%% ===== DEUTSCH =====
\recipient{ilert GmbH}{Personalabteilung}{}

\opening{Sehr geehrtes Hiring-Team,}

\vspace{8pt}\noindent\textbf{1. Warum ilert GmbH?}\par\vspace{3pt}

Die strategische Ausrichtung von ilert auf eine KI-zentrierte autonome Schicht für das Incident Response Management bietet eine spannende Gelegenheit, mein Fachwissen in verteilten Systemen einzubringen. Ich möchte Ihr Team dabei unterstützen, eine Plattform zu skalieren, auf der KI Probleme aktiv analysiert und so den manuellen Aufwand für DevOps-Ingenieure deutlich reduziert. Mit meiner Erfahrung in React, TypeScript und Node.js schlage ich die Brücke zwischen Frontend- und Backend-Anforderungen, während meine Arbeit mit Kafka und verteilten Pipelines nahtlos zu Ihrer Datenarchitektur passt. Meine Builder-Mentalität zeigt sich unter anderem in Projekten wie IRS, in dem ich Workflows mittels FastMCP automatisiert und ein Information-Retrieval-System mit hoher Datenintegrität aufgebaut habe, was ein rigoroses Systemdesign erforderte.

\vspace{8pt}\noindent\textbf{2. Warum ich?}\par\vspace{3pt}

\bullets

  \item Verantwortung für die IRS-Plattform von A bis Z übernommen: Entwicklung von Go-Mikrodiensten mit gRPC, Aufbau eines ereignisgesteuerten Systems sowie Deployment auf Kubernetes mit Cilium.

  \item Integration mehrerer Drittanbieter über ein Adapter-Pattern, um einen zuverlässigen Datenfluss zwischen externen Anbietern und der zentralen Informations-Engine zu gewährleisten.

  \item Einsatz von Hatchet zur Workflow-Verwaltung sowie Webhooks zur Automatisierung,

\bulletsend

Ich studiere derzeit an der TU Ilmenau, mein Pflichtstudium ist abgeschlossen und ich beginne mit meiner Masterarbeit, was mir eine volle Verfügbarkeit sowie einen sofortigen Umzug nach Köln für das hybride Arbeitsmodell ermöglicht.

\closing{Mit freundlichen Grüßen,}{Radheem Bin Razi}

\end{document}
